\appendix

\section{Experimental Details}
\label{app:experimental-details}

This appendix reports implementation details recorded in the project scripts and output files. The description is restricted to the final residual-aggregation experiments implemented in \texttt{run\_coverage\_ladder.py}, \texttt{run\_residual\_placebos.py}, \texttt{run\_multiseed\_placebos.py}, and the shared utilities in \texttt{run\_compositional\_feasibility.py}. Details from earlier exploratory experiments are not repeated unless they are used by the final pipeline.

\subsection{Rolling-origin evaluation}
\label{app:rolling-origin}

The final experiments use a rolling-origin evaluation with an expanding training window. The monthly observations are ordered chronologically, and the first 48 months are used as the initial training period. All subsequent months are used as one-step-ahead evaluation origins. Given the sample period from January 2018 to March 2026, this produces 51 evaluation months. At each origin, the model is trained only on observations dated before the target month and is evaluated on the selected brand series in the target month.

The forecast horizon is one month. The code does not implement a fixed holdout test set separate from the rolling-origin procedure. Instead, each evaluation month acts as one test origin, and the training sample expands as the origin moves forward. Within each rolling-origin fold, the LightGBM wrapper separates the data available at that origin into a fitting subset and an internal validation subset. The validation subset is recalculated at each origin and consists of the last 12 months available within the current training window. This validation subset is used for early stopping.

For each value of \(K \in \{5,8,10,12,14\}\), the selected brands are determined once using the first 48 months. The ranking criterion implemented in the code is the average official market share over these initial months. The same selected brands are then used throughout the rolling-origin evaluation. This fixed Top-\(K\) design avoids the use of future market-share information and keeps the target set constant across evaluation origins, so that changes in forecast accuracy are not confounded with changes in the set of evaluated brands. The baseline system, denoted DirectK, trains a global model on the \(K\) selected brand series. The residual-aggregation system trains the same global model on the \(K\) selected brand series plus the residual series. In both systems, forecast accuracy is evaluated only on the selected named brands; the residual series is not included in the evaluation target.

\subsection{Feature construction}
\label{app:feature-construction}

The final residual-aggregation experiments use the function \texttt{component\_sales\_features}. The target variable is monthly component sales, where a component is either a selected brand or an auxiliary residual/placebo series. For a component \(j\) at month \(t\), the constructed predictors are lagged component sales at lags 1, 2, 3, 6, and 12, rolling means over 3 and 6 months, two seasonal terms, and component identity indicators.

The lagged variables are constructed separately for each component. The rolling means are also component-specific and are computed after shifting the series by one month, so they use only past observations. The seasonal terms are
\[
\sin(2\pi m_t/12), \qquad \cos(2\pi m_t/12),
\]
where \(m_t\) is the calendar month. Component identity is implemented through one-hot dummy variables generated by \texttt{pandas.get\_dummies}. The scripts therefore use explicit component dummies rather than LightGBM's native categorical-feature handling.

The official market-size series is used to construct the residual category,
\[
R_t = M_t - \sum_{i \in B_K} y_{i,t},
\]
where \(M_t\) is official market size and \(B_K\) is the selected Top-\(K\) brand set. In the final Top-\(K\) residual-aggregation experiments, market size is not used as a direct predictor in the LightGBM feature matrix. Macroeconomic variables are also not used as predictors in the final residual-aggregation experiments. Although earlier project scripts constructed lagged macroeconomic and market-size predictors for the market-structure experiments, those feature blocks are not part of the final Top-\(K\) residual-aggregation pipeline.

Before model fitting, the code drops rows with missing values in the required feature columns and target column. Predictions are truncated at zero using \texttt{np.maximum(prediction, 0.0)}. The residual-aggregation script checks that the constructed residual series is strictly positive for each \(K\); if a non-positive residual is found, the script raises an error.

\subsection{Model configuration}
\label{app:model-configuration}

All final Top-\(K\) experiments use the same pooled LightGBM regression wrapper, \texttt{lgb\_predict}. The same model configuration is used for DirectK and DirectK+Residual; the only design change is the set of series included in the pooled training panel. The LightGBM parameters recorded in the implementation are:

\begin{verbatim}
objective          = regression_l1
metric             = mae
learning_rate      = 0.03
num_leaves         = 7
max_depth          = 3
min_data_in_leaf   = 20
lambda_l1          = 0.1
lambda_l2          = 1.0
feature_fraction   = 0.9
bagging_fraction   = 0.9
bagging_freq       = 1
seed               = 20260623
deterministic      = True
force_col_wise     = True
num_threads        = 1
verbosity          = -1
\end{verbatim}

The number of boosting rounds is set to 1200, with early stopping after 60 rounds without improvement on the internal validation subset. The objective function is least absolute deviations through LightGBM's \texttt{regression\_l1} objective, and the validation metric is MAE. The implementation records no Optuna search space, no automated hyperparameter search, and no data-driven hyperparameter selection for the final residual-aggregation experiments. The hyperparameters are fixed in the source code.

\subsection{Experimental settings}
\label{app:experimental-settings}

The primary Top-\(K\) ladder uses \(K = 5, 8, 10, 12,\) and 14. For each \(K\), the DirectK and DirectK+Residual systems are evaluated on the same selected target brands and the same rolling-origin months. The comparison statistic is computed from paired monthly MAE differences. In \texttt{run\_coverage\_ladder.py}, bootstrap confidence intervals are generated by resampling these paired monthly differences with replacement. The bootstrap uses 10,000 replications and seed \(20260623 + K\). The stored comparison output reports the number of evaluation months, the mean MAE difference, the relative improvement, the share of months in which the residual-aggregation system is better, the 2.5\% and 97.5\% bootstrap quantiles, and the bootstrap probability that the residual-aggregation system is better.

The placebo mechanism experiments compare the authentic residual with time-shuffled, synthetic, and duplicated auxiliary series. The time-shuffled placebo permutes the training-period residual values within each rolling-origin fold. The synthetic placebo is generated as an AR(1)-type series with mean, variance, and estimated autocorrelation based on the available training-period residual values, with the estimated autoregressive coefficient clipped to the interval \([-0.9,0.9]\). The duplicate auxiliary series is constructed from the selected brand whose training-period mean sales is the median among the selected brands. Placebo series are regenerated inside each rolling-origin fold using only information available before the test month.

The multi-seed placebo validation is implemented in \texttt{run\_multiseed\_placebos.py}. It uses 20 independent placebo replications for each \(K\) and evaluates the time-shuffled and synthetic placebo conditions. The seed passed to placebo generation is
\[
20260623 + 100000r + 1000K + f,
\]
where \(r\) is the replicate index and \(f\) is the rolling-origin fold index. The multi-seed output records the mean placebo-minus-real MAE, median placebo-minus-real MAE, 5\% and 95\% seed quantiles, the proportion of seeds in which the authentic residual is better, the average share of months in which the authentic residual is better, and the mean placebo-minus-DirectK MAE.

The recorded implementation environment uses Python 3.12.13. The project files record the use of \texttt{pandas}, \texttt{numpy}, and \texttt{lightgbm}. The installed runtime queried during appendix preparation reported \texttt{pandas} 3.0.1 and \texttt{numpy} 2.3.5. The local dependency metadata records LightGBM 4.6.0. Versions for any additional transitive dependencies are not recorded in the experiment scripts or result files.
